\documentclass[10pt,a4paper]{amsart}
\usepackage[margin=19mm]{geometry}
\usepackage{amsmath,amssymb,mathtools}
\usepackage[scr=rsfs,cal=euler]{mathalfa}
\usepackage{xfrac}
\usepackage[unicode,hidelinks,bookmarks=false]{hyperref}
\newcommand{\ie}{i.e.\ }
\newcommand{\cf}{cf.\ }
\newcommand{\resp}{respectively\ }
\newcommand{\Subsection}{\subsection}
\providecommand{\nobreakdash}{\nobreak\hbox{-}}
\setlength{\emergencystretch}{2em}
\multlinegap0pt
\usepackage{amssymb}
\usepackage{bm}
\newtheorem{thm}{Theorem}
\newcommand{\Impart}{\mathrm{Im}}
\title{Strong-disorder expansion of the root-averaged density of states for the Anderson model on the Bethe lattice}
\author{Masahiro Kaminaga}
\date{Source: arXiv:2605.00478v1 (CC BY 4.0)}
\begin{document}
\maketitle

\noindent\emph{Source and licence.} Strong-disorder expansion of the root-averaged density of states for the Anderson model on the Bethe lattice. Version arXiv:2605.00478v1 (CC BY 4.0), distributed by arXiv under the Creative Commons Attribution 4.0 International licence, http://creativecommons.org/licenses/by/4.0/, as recorded in the arXiv licence metadata for this exact version and shown on the abstract page https://arxiv.org/abs/2605.00478v1. Full licence notice: \texttt{licenses/arxiv-2605.00478-CC-BY-4.0.txt}.\par\smallskip

\noindent\textit{Editorial scope.} Counted unit: the article's thm thm:main-dos (source0.tex lines 389--430). The article's complete original proof of it is delivered with it (source0.tex lines 1113--1187, one \texttt{proof} environment of 75 lines, which the article places away from the statement). No mathematics is written, repaired or re-derived: the statement and the proof are the article's own lines, in the article's own notation, and no retained line is substituted or re-flowed.  Retained uncounted as the source-local context the proof needs: lines 309--354; lines 359--363; lines 369--374; lines 1099--1103.  Nothing was omitted from any retained span, so the package carries no omission record.  No retained line contains a citation, so the wrapper carries no bibliography and no bibliography entry.  No retained line contains a figure environment, an image-inclusion command or drawing code, so no image is delivered and the package declares no asset dependency.  Editorial formatting only, in addition: an AMS \texttt{amsart} preamble that restates the article's own declarations and macros used by the retained lines, namely the environments \texttt{thm} on separate counters, together with the standard packages those lines need.  The retained environments print in this document as Theorem 1 in order of appearance, whereas the article prints the same items with the numbering of its own full document.  The complete native source of the article as distributed in the arXiv e-print for this exact version is bundled unchanged as \texttt{sources/199565/source0.tex}.
\begin{thm}\label{thm:main-resolvent}
Assume the local analyticity condition above. 
Fix $0<\delta<\delta_0$.
Then, there exist $\lambda_0>0$ and functions $M_n$, $n=0,1,2,\ldots$,
holomorphic on $\Omega_\delta(I)$ such that, 
for every $\lambda\ge\lambda_0$, the function
$$
\zeta\mapsto m_\lambda(\lambda\zeta),
\qquad \Impart\zeta>0,
$$
extends holomorphically from $\mathbb C_+$ to $\Omega_\delta(I)$, and
the series
$$
\sum_{n=0}^{\infty}\lambda^{-n-1}M_n(\zeta)
$$
converges uniformly on $\Omega_\delta(I)$ to this holomorphic continuation.

Moreover, for every $N\ge0$, there exists $C_{N,\delta}>0$ such that
\begin{equation}
\label{eq:main-resolvent-expansion-zeta}
m_\lambda(\lambda\zeta)=\sum_{n=0}^{N}\lambda^{-n-1}M_n(\zeta)
+\widetilde R_{N,\lambda}(\zeta),\qquad\zeta\in\Omega_\delta(I),
\end{equation}
with the uniform bound
$$
|\widetilde R_{N,\lambda}(\zeta)|
\le
C_{N,\delta}\lambda^{-N-2}.
$$
The coefficient functions are given by
\begin{equation}
\label{eq:Mn-formula-statement}
M_n(\zeta)
=
(-1)^n
\sum_{\gamma\in\Gamma_n(0,0)}
\prod_{x:\,\nu_\gamma(x)>0}
s_{\nu_\gamma(x)}(\zeta),
\qquad
\zeta\in\Omega_\delta(I).
\end{equation}
In particular, $M_n\equiv0$ for every odd $n$, and
$$
M_0(\zeta)=s_1(\zeta).
$$
\end{thm}

\begin{equation}
\label{eq:main-resolvent-expansion}
m_\lambda(z)=\sum_{n=0}^{N}\lambda^{-n-1}M_n(z/\lambda)+R_{N,\lambda}(z),
\qquad z/\lambda\in\Omega_\delta(I),
\end{equation}

\begin{equation}
\label{eq:main-resolvent-remainder}
|R_{N,\lambda}(z)|
\le
C_{N,\delta}\lambda^{-N-2}.
\end{equation}

\begin{equation}
\label{eq:dos-stieltjes}
m_\lambda(z)=\int_{\mathbb R}\frac{dN_\lambda(E)}{E-z},
\qquad\Impart z>0,
\end{equation}

\begin{thm}
\label{thm:main-dos}
Assume the same hypotheses as in Theorem~\ref{thm:main-resolvent}. Fix
$0<\delta<\delta_0$ and $N\ge0$, and let $\lambda_0$ and
$C_{N,\delta}$ be the constants in Theorem~\ref{thm:main-resolvent}.
Then, for every $\lambda\ge\lambda_0$, the density of states measure of
$H_{\lambda,\omega}$ is absolutely continuous on $\lambda I$, and its
density $n_\lambda(E)$ is real analytic on $\lambda I$. Moreover, for
$\lambda\ge\lambda_0$ and $\xi\in I$,
\begin{equation}
\label{eq:main-dos-expansion}
n_\lambda(\lambda\xi)
=
\sum_{n=0}^{N}\lambda^{-n-1}a_n(\xi)
+
r_{N,\lambda}(\xi),
\end{equation}
where
\begin{equation}
\label{eq:an-def}
a_n(\xi)
=
\frac{1}{\pi}\Impart M_n(\xi),
\qquad
\xi\in I.
\end{equation}
Here $M_n(\xi)$ denotes the value at $\xi\in I$ of the holomorphic
continuation of $M_n$ obtained in Theorem~\ref{thm:main-resolvent}.
The remainder satisfies
\begin{equation}
\label{eq:rn-bound}
|r_{N,\lambda}(\xi)|
\le
\frac{1}{\pi}C_{N,\delta}\lambda^{-N-2}.
\end{equation}
In particular,
\begin{equation}
\label{eq:a0-rho}
a_0(\xi)=\rho(\xi),
\end{equation}
and $a_n\equiv0$ for every odd $n$.
\end{thm}

\begin{proof}
We prove Theorem~\ref{thm:main-dos}.
By Theorem~\ref{thm:main-resolvent}, for every sufficiently large $\lambda$ the function
$\zeta\mapsto m_\lambda(\lambda\zeta)$ is holomorphic on $\Omega_\delta(I)$.
Hence, the function
$$
E\mapsto m_\lambda(E)
$$
is holomorphic on the neighborhood $\lambda\Omega_\delta(I)$ of $\lambda I$.
Thus, its upper half-plane boundary value on $\lambda I$
exists and is equal to the restriction of this holomorphic continuation.

By the Stieltjes inversion formula applied to \eqref{eq:dos-stieltjes}, 
with the sign convention used there, the
density corresponding to the upper half-plane boundary value is
$(1/\pi)\Impart m_\lambda(E+i0)$. 
By the Stieltjes inversion formula applied to
\eqref{eq:dos-stieltjes}, with the sign convention used there, for every
$\varphi\in C_0^\infty(\lambda I)$ one has
$$
\int_{\lambda I}\varphi(E)\,dN_\lambda(E)
=
\lim_{\varepsilon\downarrow0}\frac{1}{\pi}\int_{\lambda I}
\varphi(E)\Impart m_\lambda(E+i\varepsilon)\,dE .
$$
Since the upper half-plane branch of $m_\lambda$ extends holomorphically
across $\lambda I$, this convergence is locally uniform on $\lambda I$.
Hence, we have
$$
\int_{\lambda I}\varphi(E)\,dN_\lambda(E)
=
\int_{\lambda I}
\varphi(E)\frac{1}{\pi}\Impart m_\lambda(E)dE.
$$
Therefore, $dN_\lambda$ is absolutely continuous on $\lambda I$, 
and its density is given by
\begin{equation}
\label{eq:density-boundary}
n_\lambda(E)=\frac{1}{\pi}\Impart m_\lambda(E),\qquad E\in\lambda I,
\end{equation}
where $m_\lambda(E)$ denotes the value at $E\in\lambda I$ of the holomorphic continuation. 
Since $m_\lambda(E)$ is real analytic on $\lambda I$, the same is true for $n_\lambda(E)$.

Now let $\xi\in I$. Substituting $z=\lambda\xi$ into \eqref{eq:main-resolvent-expansion} gives
\begin{equation}
\label{eq:boundary-at-xi}
m_\lambda(\lambda\xi)=\sum_{n=0}^{N}\lambda^{-n-1}M_n(\xi)
+R_{N,\lambda}(\lambda\xi).
\end{equation}
Here $M_n(\xi)$ denotes the value at $\xi\in I$ of the holomorphic
continuation obtained in Theorem~\ref{thm:main-resolvent}. 
Taking the imaginary part and dividing by $\pi$, we obtain
\eqref{eq:main-dos-expansion} with $a_n$ defined by \eqref{eq:an-def} and with
$$
r_{N,\lambda}(\xi)=\frac{1}{\pi}\Impart R_{N,\lambda}(\lambda\xi).
$$
The bound \eqref{eq:rn-bound} follows from \eqref{eq:main-resolvent-remainder}.

It remains to identify $a_0$. 
Since $M_0=s_1$, we have
$$
a_0(\xi)=\frac{1}{\pi}\Impart s_1(\xi),
$$
where $s_1(\xi)$ denotes the value at $\xi\in I$ of the holomorphic
continuation of the single-site Stieltjes transform. 
Since $s_1$ is the Stieltjes transform of $\mu$ and $\mu$ has density $\rho$ on $I$, 
the Stieltjes inversion formula gives
$$
\frac{1}{\pi}\Impart s_1(\xi)=\rho(\xi)
$$
for $\xi\in I$. 
This proves \eqref{eq:a0-rho}. 
The vanishing of odd $a_n$ follows from the vanishing of odd $M_n$ in Theorem~\ref{thm:main-resolvent}.
This completes the proof.
\end{proof}

\end{document}
